% ============ Appendix C. The exterior algebra (原书 pp.283-286) ============
\chapter*{Appendix C}
\addcontentsline{toc}{chapter}{Appendix C: The exterior algebra}
\markboth{Appendix C}{The exterior algebra}
\renewcommand{\thetheorem}{C\arabic{theorem}}
\setcounter{theorem}{0}

\vspace*{-0.5cm}
\begin{center}
  {\Large\bfseries The exterior algebra}
\end{center}
\vspace{0.3cm}

\begin{enumerate}[(1)]
  \item Let $M$ and $N$ be modules over a ring $A$. An $r$-multilinear map
        $\varphi\colon M^r = M \times \cdots \times M \to N$ from the direct product of $r$ copies of
        $M$ is said to be \emph{alternating} if $\varphi(x_1, \ldots, x_r) = 0$ whenever any of the
        elements $x_1, \ldots, x_r$ appears more than once. If $\varphi$ is alternating then for
        any $x_1, \ldots, x_r \in M$ we have
        \[
          \varphi(x_1, \ldots, x_{i-1}, x_i + x_j, x_{i+1}, \ldots, x_{j-1}, x_i + x_j,
              x_{j+1}, \ldots) = 0,
        \]
        and expanding out the left-hand side gives
        \[
          \varphi(x_1, \ldots, x_i, \ldots, x_j, \ldots) + \varphi(x_1, \ldots, x_j, \ldots,
              x_i, \ldots) = 0.
        \]
        In other words, on interchanging two of its entries, $\varphi$ changes sign.

        The \emph{$r$th exterior product} of $M$ is defined as the module $N_r$ having a
        universal alternating $r$-fold multilinear map $f_0\colon M^r \to N_r$, that is a map
        satisfying the property that every alternating $r$-fold multilinear map
        $f\colon M^r \to N$ factorises as $f = h \circ f_0$ for a unique $A$-linear map
        $h\colon N_r \to N$. We write $N_r = \bigwedge^r M$, and use
        $x_1 \wedge \cdots \wedge x_r$ to denote $f_0(x_1, \ldots, x_r)$. To prove the existence of the
        exterior product, let $N_r$ be the quotient of the $r$-fold tensor product
        $M \otimes \cdots \otimes M$ by the submodule generated by elements of the form
        $x_1 \otimes \cdots \otimes x \otimes \cdots \otimes x \otimes \cdots \otimes x_r$. Then $N_r$ satisfies the above
        condition. The fact that the exterior product is uniquely determined up to
        isomorphism is obvious from the definition.

  \item If $M$ is a free $A$-module of rank $n$, with basis $e_1, \ldots, e_n$, then
        $\bigwedge^r M$ is zero if $r > n$, and if $r \le n$ is the free $A$-module of rank
        $\binom{n}{r}$ with basis $(e_{i_1} \wedge \cdots \wedge e_{i_r} \mid 1 \le i_1 < \cdots < i_r \le n)$.
        (If $r > n$ this is easy; if $r \le n$ then the elements given above obviously
        generate $\bigwedge^r M$, and the fact that they are linearly independent can
        also be proved by reducing to the theory of determinants.)

  \item However, if $I \subset A$ is an ideal, then $\bigwedge^r A = 0$, but nevertheless
        $\bigwedge^r I$ is not necessarily $0$. For example, let $k$ be a field, $x$ and $y$
        indeterminates, and $A = k[x, y]$; if $I = xA + yA$ then $\bigwedge^2 I \ne 0$. Indeed,
        we can define $\varphi\colon I \times I \to k = A/I$ by
        \[
          \varphi(ax + by, a'x + b'y) = ab' - a'b,
        \]
        and it is then easy to check that $\varphi$ is alternating and bilinear, with
        $\varphi(x, y) = 1$, so that $\varphi \ne 0$, and we must have $\bigwedge^2 I \ne 0$.

  \item The operation of taking the exterior product commutes with extensions
        of scalars. That is, let $B$ be an $A$-algebra and $M$ an $A$-module, and set
        $M \otimes_A B = M_B$. Then $(\bigwedge^r M) \otimes_A B = \bigwedge^r M_B$, where of course
        $\bigwedge^r$ on the right-hand side refers to the exterior product of $B$-modules.
        For the proof, according to Appendix A, Formula~11, we have
        \[
          M_B \otimes_B \cdots \otimes_B M_B = (M \otimes_A \cdots \otimes_A M) \otimes_A B,
        \]
        so that if we let $f_0$ be the composite
        \[
          (M_B)^r = \underbrace{M \otimes_A \cdots \otimes_A M}_{r} \otimes_A B \to
              \Big(\bigwedge^r M\Big) \otimes_A B,
        \]
        then $f_0$ is an alternating $r$-fold $B$-multilinear map. Write
        $v\colon M \to M_B$ for the natural map $x \mapsto x \otimes 1$. Let $N$ be a $B$-module and
        $\varphi\colon (M_B)^r \to N$ be an alternating $r$-fold $B$-multilinear map. Then
        $\varphi$ induces an alternating $r$-fold $A$-multilinear map $\varphi\colon M^r \to N$,
        and therefore an $A$-linear map $\bigwedge^r M \to N$, and finally a $B$-linear map
        $(\bigwedge^r M) \otimes_A B \to N$ which we denote $h$. Then on $v(M^r)$ the two maps
        $\varphi$ and $h \circ f_0$ coincide; but $M_B$ is generated over $B$ by $v(M)$, so that
        $\varphi = h \circ f_0$. Thus $f_0$ has the universal property, and we can think of
        $(\bigwedge^r M) \otimes_A B$ as $\bigwedge^r M_B$.
\end{enumerate}

\begin{theorem}\label{thm:C1}
Let $A$ be an integral domain with field of fractions $K$, and let
$I_1, \ldots, I_r$ be ideals of $A$. Set $M = I_1 \otimes \cdots \otimes I_r$, and let $T$ be the torsion
submodule of $\bigwedge^r M$; then $(\bigwedge^r M)/T \cong I_1 \cdots I_r$. Therefore if
$J_1, \ldots, J_r$ are ideals of $A$ such that $I_1 \otimes \cdots \otimes I_r \cong J_1 \otimes \cdots \otimes J_r$
we have $I_1 \cdots I_r = J_1 \cdots J_r$.
\end{theorem}

\begin{proof}
We have $M_K \cong K \otimes \cdots \otimes K$ (the direct sum of $r$ copies of $K$), so that
$(\bigwedge^r M) \otimes K = \bigwedge^r M_K \cong K$. The kernel of the natural map
$\bigwedge^r M \to (\bigwedge^r M) \otimes K \cong K$ is obviously $T$ (since tensoring with $K$ is
the same thing as the localisation with respect to the zero ideal of $A$, see
\S\ref{sec:4}). In addition, the image is $I_1 \cdots I_r$. Indeed, viewing each $I_i$ as a
submodule of $K$, we can assume that the map is
\[
  \bigwedge^r (I_1 \otimes \cdots \otimes I_r) \to K,
\]
and since $\xi_i = \sum_{j=1}^{r} a_{ij}e_j \in K^r$ for $\xi_i \in I_i \otimes K$ we have
$\xi_1 \wedge \cdots \wedge \xi_r = \det(a_{ij})e_1 \wedge \cdots \wedge e_r$; it is clear from this that the above
map has image $I_1 \cdots I_r e_1 \wedge \cdots \wedge e_r$.
\end{proof}

If $A$ is a Dedekind ring then it is known that $I_1 \otimes \cdots \otimes I_r \cong A^{r-1} \otimes
I_1 \cdots I_r$ (see for example [B7], \S4, Prop.~24).

\begin{theorem}\label{thm:C2}
Let $A$ be a ring and $M$, $N$ $A$-modules. Then
\[
  \bigwedge^r (M \oplus N) = \bigoplus_{s + t = r} \Big[\Big(\bigwedge^s M\Big) \otimes_A
      \Big(\bigwedge^t N\Big)\Big].
\]
\end{theorem}

\begin{proof}
$\bigotimes^r(M \oplus N)$ can be written as a direct sum of all possible $r$-fold tensor
products of copies of $M$ and $N$ (with $2^r$ summands). Write $L_{s,t}$ for the
submodule which is the direct sum of all tensor products involving $s$ copies of
$M$ and $t$ copies of $N$ (with $s + t = r$). Thus
\[
  \bigotimes^r(M \oplus N) = \bigoplus_{s + t = r} L_{s,t}.
\]
For example, when $r = 2$ we have $L_{2,0} = M \otimes M$,
$L_{1,1} = (M \otimes N) \oplus (N \otimes M)$ and $L_{0,2} = N \otimes N$. Now let $Q$ be the
submodule of $\bigotimes^r(M \oplus N)$ generated by all elements of the form
$\cdots \otimes x \otimes \cdots \otimes x \otimes \cdots$; we have
$Q = \bigoplus[Q \cap L_{s,t}]$. We see at once that $Q \cap L_{s,t}$ is the submodule generated by
all elements of the forms $\cdots \otimes y \otimes \cdots \otimes y \otimes \cdots$ (with either $y \in M$
or $y \in N$), and $x \otimes y \otimes \xi \otimes y + \xi \otimes y \otimes x \otimes y$ (with $x \in M$ and
$y \in N$). Thus one sees easily that
\[
  \bigwedge^r (M \oplus N) = \Big(\bigotimes^r(M \oplus N)\Big)/Q = \bigoplus_{s + t = r}
      (L_{s,t}/Q \cap L_{s,t}),
\]
and
\[
  L_{s,t}/L_{s,t} \cap Q \cong \Big(\bigwedge^s M\Big) \otimes \Big(\bigwedge^t N\Big).
\]
\end{proof}

\begin{enumerate}[(5)]
  \item Let $A$ be a commutative ring. We say that a (possibly non-commutative)
        $A$-algebra $E$ is a \emph{skew-commutative graded algebra} if it has a direct sum
        decomposition $E = \bigoplus_{n \ge 0} E_n$ as an $A$-module, such that
        \begin{enumerate}[(i)]
          \item $E_pE_q = E_{p+q}$;
          \item $xy = (-1)^{pq}yx$ for $x \in E_p$ and $y \in E_q$;
          \item $x^2 = 0$ for $x \in E_{2k+1}$.
        \end{enumerate}
        For such an algebra $E$, a \emph{skew-derivation} is an $A$-linear map $d\colon E \to E$
        such that
        \begin{enumerate}[(a)]
          \item $d(E_n) \subset E_{n-1}$;
          \item[($\beta$)] $d(xy) = (dx)y + (-1)^px(dy)$ for $x \in E_p$, $y \in E_q$.
        \end{enumerate}
  \item Let $A$ be a ring and $M$ an $A$-module. We show how to define an
        $A$-bilinear map $\psi\colon (\bigwedge^p M) \times (\bigwedge^q M) \to \bigwedge^{p+q} M$. If we
        define $\varphi\colon M^p \times M^q \to \bigwedge^{p+q} M$ by
        $\varphi(x_1, \ldots, x_p, y_1, \ldots, y_q) = x_1 \wedge \cdots \wedge x_p \wedge y_1 \wedge \cdots \wedge y_q$,
        then for fixed $y_1, \ldots, y_q$ this is an alternating $p$-multilinear map from
        $M^p$ to $\bigwedge^{p+q} M$, so that there is a map $\Phi\colon (\bigwedge^p M) \times M^q \to
        \bigwedge^{p+q} M$ such that $\Phi(\xi, y_1, \ldots, y_q)$ is linear in $\xi$ and satisfies
        $\Phi(\xi, y_1, \ldots, y_q) = x_1 \wedge \cdots \wedge x_p \wedge y_1 \wedge \cdots \wedge y_q$ if
        $\xi = x_1 \wedge \cdots \wedge x_p$. Now for fixed $\xi$, $\Phi$ is alternating and multilinear in
        $y_1, \ldots, y_q$, defining a bilinear map $\psi\colon (\bigwedge^p M) \times (\bigwedge^q M) \to
        \bigwedge^{p+q} M$ such that $\psi(\xi, y_1 \wedge \cdots \wedge y_q) = \Phi(\xi, y_1, \ldots, y_q)$. For
        $\xi \in \bigwedge^p M$ and $\eta \in \bigwedge^q M$ we write $\xi \wedge \eta$ for $\psi(\xi, \eta)$; if
        $\xi = \sum_\alpha x_\alpha^{(1)} \wedge \cdots \wedge x_\alpha^{(p)}$ and
        $\eta = \sum_\beta y_\beta^{(1)} \wedge \cdots \wedge y_\beta^{(q)}$ then
        $\xi \wedge \eta = \sum_{\alpha, \beta} x_\alpha^{(1)} \wedge \cdots \wedge x_\alpha^{(p)}
        \wedge y_\beta^{(1)} \wedge \cdots \wedge y_\beta^{(q)}$. (It might be tempting to make the
        definition directly in terms of this formula, but the expression of $\xi$ and
        $\eta$ in the above form is non-unique, so that this requires an awkward
        proof.) The multiplication $\wedge$ satisfies the associative law
        $\xi \wedge (\eta \wedge \zeta) = (\xi \wedge \eta) \wedge \zeta$, so that we can use it to define a
        product on $\bigoplus_{n \ge 0} \bigwedge^n M$ (we set $\bigwedge^0 M = A$), which becomes an
        $A$-algebra, and it is easy to see that this satisfies the conditions for a
        skew-commutative algebra. We write $\bigwedge M$ for this algebra, and call it
        the \emph{exterior algebra} of $M$.

        Given any linear map $\alpha\colon M \to A$, there exists a unique skew-derivation
        $d$ of $\bigwedge M$ such that $d$ coincides with $\alpha$ on
        $\bigwedge^1 M = M$. The uniqueness is clear from the fact that $\bigwedge M$ is
        generated as an $A$-algebra by $M$: we must have
        \[
          d(x_1 \wedge \cdots \wedge x_p) = \sum_{i=1}^{p} (-1)^{i-1}\alpha(x_i)x_1
              \wedge \cdots \wedge \widehat{x_i} \wedge \cdots \wedge x_p.
        \]
        Conversely, the existence follows easily from the fact that the right-hand
        side defines an alternating $p$-fold multilinear map of $x_1, \ldots, x_p$.

        In particular, let $M$ be a free $A$-module of rank $n$ with basis
        $e_1, \ldots, e_n$, so that $M = Ae_1 \oplus \cdots \oplus Ae_n$. Then taking arbitrary
        elements $c_1, \ldots, c_n \in A$, we can define $\alpha\colon M \to A$ by
        $\alpha(e_i) = c_i$, and the skew-derivation of $\bigwedge M$ takes the form
        \[
          d(e_{i_1} \wedge \cdots \wedge e_{i_p}) = \sum_{j=1}^{p} (-1)^{j-1}c_{i_j}
              e_{i_1} \wedge \cdots \wedge \widehat{e_{i_j}} \wedge \cdots \wedge e_{i_p}.
        \]
        This can be identified with the differential operator of the Koszul complex
        $K_{\bullet}^{c_1, \ldots, c_n}$ discussed in \S\ref{sec:16}. Thus the Koszul complex can be
        thought of as the exterior algebra $\bigwedge(Ae_1 \oplus \cdots \oplus Ae_n)$ with the
        skew-derivation defined by $d(e_i) = c_i$.
\end{enumerate}

% ============ 习题 ============
\section*{Exercise to Appendix C}

\noindent\textbf{C.1.}\quad Let $(A, m, k)$ be a local ring and $M$ a finitely generated
$A$-module. Prove that $\min\{r \mid \bigwedge^r M \ne 0\}$ is equal to the minimal number of
generators of $M$.
